% Options for packages loaded elsewhere
\PassOptionsToPackage{unicode}{hyperref}
\PassOptionsToPackage{hyphens}{url}
\documentclass[
  11pt,
]{article}
\usepackage{xcolor}
\usepackage[left=3cm,right=5cm,top=2cm,bottom=2cm]{geometry}
\usepackage{amsmath,amssymb}
\setcounter{secnumdepth}{5}
\usepackage{iftex}
\ifPDFTeX
  \usepackage[T1]{fontenc}
  \usepackage[utf8]{inputenc}
  \usepackage{textcomp} % provide euro and other symbols
\else % if luatex or xetex
  \usepackage{unicode-math} % this also loads fontspec
  \defaultfontfeatures{Scale=MatchLowercase}
  \defaultfontfeatures[\rmfamily]{Ligatures=TeX,Scale=1}
\fi
\usepackage{lmodern}
\ifPDFTeX\else
  % xetex/luatex font selection
\fi
% Use upquote if available, for straight quotes in verbatim environments
\IfFileExists{upquote.sty}{\usepackage{upquote}}{}
\IfFileExists{microtype.sty}{% use microtype if available
  \usepackage[]{microtype}
  \UseMicrotypeSet[protrusion]{basicmath} % disable protrusion for tt fonts
}{}
\makeatletter
\@ifundefined{KOMAClassName}{% if non-KOMA class
  \IfFileExists{parskip.sty}{%
    \usepackage{parskip}
  }{% else
    \setlength{\parindent}{0pt}
    \setlength{\parskip}{6pt plus 2pt minus 1pt}}
}{% if KOMA class
  \KOMAoptions{parskip=half}}
\makeatother
\usepackage{longtable,booktabs,array}
\newcounter{none} % for unnumbered tables
\usepackage{calc} % for calculating minipage widths
% Correct order of tables after \paragraph or \subparagraph
\usepackage{etoolbox}
\makeatletter
\patchcmd\longtable{\par}{\if@noskipsec\mbox{}\fi\par}{}{}
\makeatother
% Allow footnotes in longtable head/foot
\IfFileExists{footnotehyper.sty}{\usepackage{footnotehyper}}{\usepackage{footnote}}
\makesavenoteenv{longtable}
\usepackage{graphicx}
\makeatletter
\newsavebox\pandoc@box
\newcommand*\pandocbounded[1]{% scales image to fit in text height/width
  \sbox\pandoc@box{#1}%
  \Gscale@div\@tempa{\textheight}{\dimexpr\ht\pandoc@box+\dp\pandoc@box\relax}%
  \Gscale@div\@tempb{\linewidth}{\wd\pandoc@box}%
  \ifdim\@tempb\p@<\@tempa\p@\let\@tempa\@tempb\fi% select the smaller of both
  \ifdim\@tempa\p@<\p@\scalebox{\@tempa}{\usebox\pandoc@box}%
  \else\usebox{\pandoc@box}%
  \fi%
}
% Set default figure placement to htbp
\def\fps@figure{htbp}
\makeatother
% definitions for citeproc citations
\NewDocumentCommand\citeproctext{}{}
\NewDocumentCommand\citeproc{mm}{%
  \begingroup\def\citeproctext{#2}\cite{#1}\endgroup}
\makeatletter
 % allow citations to break across lines
 \let\@cite@ofmt\@firstofone
 % avoid brackets around text for \cite:
 \def\@biblabel#1{}
 \def\@cite#1#2{{#1\if@tempswa , #2\fi}}
\makeatother
\newlength{\cslhangindent}
\setlength{\cslhangindent}{1.5em}
\newlength{\csllabelwidth}
\setlength{\csllabelwidth}{3em}
\newenvironment{CSLReferences}[2] % #1 hanging-indent, #2 entry-spacing
 {\begin{list}{}{%
  \setlength{\itemindent}{0pt}
  \setlength{\leftmargin}{0pt}
  \setlength{\parsep}{0pt}
  % turn on hanging indent if param 1 is 1
  \ifodd #1
   \setlength{\leftmargin}{\cslhangindent}
   \setlength{\itemindent}{-1\cslhangindent}
  \fi
  % set entry spacing
  \setlength{\itemsep}{#2\baselineskip}}}
 {\end{list}}
\usepackage{calc}
\newcommand{\CSLBlock}[1]{\hfill\break\parbox[t]{\linewidth}{\strut\ignorespaces#1\strut}}
\newcommand{\CSLLeftMargin}[1]{\parbox[t]{\csllabelwidth}{\strut#1\strut}}
\newcommand{\CSLRightInline}[1]{\parbox[t]{\linewidth - \csllabelwidth}{\strut#1\strut}}
\newcommand{\CSLIndent}[1]{\hspace{\cslhangindent}#1}
\setlength{\emergencystretch}{3em} % prevent overfull lines
\providecommand{\tightlist}{%
  \setlength{\itemsep}{0pt}\setlength{\parskip}{0pt}}
\usepackage{amsmath}
\usepackage{amssymb}
\usepackage{lineno}
\linenumbers
\usepackage{setspace}
\setstretch{1.1}
\usepackage{bookmark}
\IfFileExists{xurl.sty}{\usepackage{xurl}}{} % add URL line breaks if available
\urlstyle{same}
\hypersetup{
  pdftitle={Test-procedure and trial-design choices for the biomarker-treatment interaction in aggregated N-of-1 trials: classical RM-ANOVA{,} mixed-effects{,} and the cycle-by-period protocol grid},
  pdfauthor={pmsimstats team},
  hidelinks,
  pdfcreator={LaTeX via pandoc}}

\title{Test-procedure and trial-design choices for the biomarker-treatment interaction in aggregated N-of-1 trials: classical RM-ANOVA, mixed-effects, and the cycle-by-period protocol grid}
\author{pmsimstats team}
\date{2026-05-21}

\begin{document}
\maketitle

{
\setcounter{tocdepth}{2}
\tableofcontents
}
\section{Abstract}\label{abstract}

\textbf{Background.} Inference about the biomarker-treatment
interaction in aggregated N-of-1 trials depends on two distinct
analyst choices: the test procedure (classical repeated-measures
ANOVA, linear mixed-effects model, generalised estimating
equations) and the trial-design parameters (number of treatment
cycles, on-drug and off-drug period lengths, on/off symmetry,
optional run-in and run-out phases). Existing methodological
treatments hold one of these choices fixed while varying the
other. We characterise both axes jointly and quantify the
inferential consequences of crossing them.

\textbf{Methods.} We derive the closed-form strict RM-ANOVA
\(F\)-statistic for the biomarker-stratified treatment-by-time
interaction in a balanced split-plot design with categorical
predictors and the sphericity assumption, and relate it to the
Wald \(t\)-test for the corresponding fixed-effects coefficient in
a linear mixed model with continuous-time AR(1) residual
correlation. We then conduct a Monte Carlo simulation study,
designed and reported under the\textsuperscript{\citeproc{ref-morris2019}{1}} ADEMP framework,
that crosses four test procedures (strict RM-ANOVA, the
standard pmsimstats linear-mixed analysis with \texttt{corCAR1}, and a
generalised estimating equations comparator with naive and with
Mancl-DeRouen small-sample sandwich variance) against the
biomarker-moderation coefficient \(c_{bm} \in \{0, 0.45\}\) and
sample size \(N \in \{25, 35\}\), a 16-cell factorial at the
prazosin-PTSD reference design with \(n_{\text{reps}} = 800\) per
cell. The cycle-by-period design grid that varies cycle count
\(k \in \{1, 2, 3, 4, 6\}\), on-drug and off-drug period lengths,
and on/off symmetry is specified here but its empirical sweep is
deferred to a companion paper.

\textbf{Results.} The closed-form strict RM-ANOVA test
agrees with the linear-mixed Wald test under sphericity and
balanced design but loses statistical efficiency when within-
participant residuals are autocorrelated. In the simulation
study reported here, the linear-mixed analysis dominates
strict RM-ANOVA on power for the biomarker-treatment
interaction at both sample sizes evaluated, with the power gap
widening from approximately 0.19 at \(N = 25\) to 0.28 at
\(N = 35\), while holding nominal or sub-nominal type I error.
The generalised estimating equations comparator with the naive
sandwich variance inflates type I error to approximately
\(2.4\times\) nominal at \(N = 25\) and \(1.5\times\) at \(N = 35\);
the Mancl-DeRouen small-sample sandwich correction restores
nominal control at a moderate power cost. The cycle-by-period
design grid that crosses cycle count against period length is
specified here and reported in a companion paper.

\textbf{Conclusions.} Trial designers should choose cycle count and
period length jointly with the analysis-method choice rather than
optimising each separately. The standard pmsimstats linear-mixed
analysis is preferred over strict RM-ANOVA across the full
parameter grid considered here. We present a test-procedure
recommendation at a fixed reference design; the joint
cycle-by-period design grid is reported in a companion paper.

\section{Introduction}\label{introduction}

\subsection{A claim about how design and test choices interact}\label{a-claim-about-how-design-and-test-choices-interact}

This paper develops the methodological case for treating the
analyst's two principal choices in an aggregated N-of-1 trial
\textbf{jointly} rather than separately. The two choices are the
\textbf{test procedure} for the biomarker-by-treatment interaction
(strict repeated-measures ANOVA, the linear mixed-effects model
with continuous-time AR(1) residual correlation, the
mixed-effects model for repeated measures, generalised
estimating equations) and the \textbf{trial-design parameters} that
govern within-subject treatment cycling (number of treatment
cycles, on-drug and off-drug period lengths, on/off symmetry,
optional run-in and run-out phases). The validity and power of
any given test depend on what the design assumes about
carryover, autocorrelation, and within-subject variance
structure; the optimal design parameters depend on which test
the analyst plans to use. Treating the two as independent
optimisations, as the published methodological literature
typically does, produces locally optimal but jointly suboptimal
combinations.

The case we develop is two-part. \textbf{First}, that a joint
design-by-test sensitivity study has not been performed for the
aggregated N-of-1 design class, and the absence of one limits
the field's ability to make defensible joint
recommendations. \textbf{Second}, that the linear mixed-effects test
with continuous-time AR(1) residual correlation, used in
conjunction with the Hendrickson hybrid open-label-plus-blinded-
discontinuation design\textsuperscript{\citeproc{ref-hendrickson2020}{2}}, dominates the strict
classical RM-ANOVA test in essentially every parameter cell
where within-participant residuals depart from compound
symmetry, which is the empirically common case in chronic-
condition trials. Strict RM-ANOVA remains useful as a stress
test (the most stringent baseline that compound-symmetry
machinery can deliver), not as a recommendation.

\subsection{Why the question matters for chronic-condition pharmacology}\label{why-the-question-matters-for-chronic-condition-pharmacology}

Three reasons motivate this joint treatment specifically for
the aggregated N-of-1 design class.

\textbf{Within-participant residuals are autocorrelated by physiology,
not by analytic convenience.} Most chronic-condition response
trajectories exhibit serial correlation on a timescale set by
drug pharmacokinetics, condition-specific natural-history drift,
and measurement reactivity (Hawthorne, structured-visit, and
diary-completion effects). Continuous-time AR(1) residual
correlation is the right structural prior for this physiology,
but it does not satisfy the sphericity assumption required by
strict RM-ANOVA.\textsuperscript{\citeproc{ref-haverkamp2017sphericity}{3}} show through
simulation that as the number of measurement occasions
increases, the rmANOVA F-statistic inflates Type I error rapidly
under sphericity violations, and that Huynh-Feldt and
Greenhouse-Geisser corrections\textsuperscript{\citeproc{ref-huynh1976}{4},\citeproc{ref-greenhouse1959}{5}}
recover nominal levels only conditionally.\textsuperscript{\citeproc{ref-blanca2023sphericity}{6}}
extend this with finer epsilon resolution and recommend
Greenhouse-Geisser when \(\epsilon < 0.6\). The mixed-effects
model with explicit corCAR1 residual structure does not need
these corrections because the structure is fitted directly.

\textbf{The biomarker-treatment interaction is harder to detect than
the average treatment effect.} The PATH framework\textsuperscript{\citeproc{ref-kent2020path}{7},\citeproc{ref-kent2020pathee}{8}} makes the methodological
distinction precise: effect-modelling (the predictive-HTE
question) requires more design-aware machinery than
risk-modelling (the average-effect question). The within-
participant contrasts that contribute to the interaction
estimator are sensitive to the trial-design parameters in ways
the main-effect analysis is not, and a test procedure that
underpowers the interaction term will systematically miss
predictive-biomarker signal even at sample sizes adequate for
the average-effect test.

\textbf{Cycle count and period length jointly determine the
within-subject contrast structure.}\textsuperscript{\citeproc{ref-senn1994crossover}{9}} and\textsuperscript{\citeproc{ref-senn2011individual}{10}} establish, in the broader cross-over
literature, that adding within-subject treatment cycles
contributes information about the patient-by-treatment
interaction variance that no cross-sectional comparator can
supply. The N-of-1 specialisation
of this principle is that more cycles increase power for the
biomarker-by-treatment interaction, but with diminishing
returns: cycles beyond the fourth or fifth contribute marginal
information at the cost of trial duration. Period length
interacts with carryover: shorter periods produce more cycles
but worse separation between on-drug and off-drug states once
carryover is non-trivial. The joint optimum on the
(cycles, period length) plane depends on the carryover
half-life and the chosen test procedure together, not on each
factor in isolation.

\subsection{What the published N-of-1 methodology literature does and does not do}\label{what-the-published-n-of-1-methodology-literature-does-and-does-not-do}

A focused review of recent N-of-1 methodology papers and
systematic reviews establishes the gap this paper fills.

\textbf{The dominant analytic models do not exploit AR(1) residual
correlation.} The systematic review of\textsuperscript{\citeproc{ref-natesanbatley2023}{11}}
documents that 83.8\% of the published N-of-1 studies they
catalogue ignore autocorrelation entirely;\textsuperscript{\citeproc{ref-stunnenberg2022}{12}} and\textsuperscript{\citeproc{ref-shaffer2018}{13}} reach similar conclusions in neurology and
behavioural-medicine subdomains. CONSORT extension reporting
standards\textsuperscript{\citeproc{ref-vohra2015cent}{14},\citeproc{ref-shamseer2015cent}{15}} codify reporting
quality but do not address the analytic-model decision in any
detail;\textsuperscript{\citeproc{ref-kwasnicka2019}{16}} identifies the analytic-method gap as
the highest-priority methodological problem in the field.

\textbf{Test-procedure comparisons are uncommon and limited in
scope.} The MMRM consensus established by Mallinckrodt and
colleagues\textsuperscript{\citeproc{ref-mallinckrodt2003}{17}--\citeproc{ref-prakash2008}{20}} applies to
parallel-group longitudinal trials with monotone missingness,
not to aggregated N-of-1 designs. The GEE methodology of\textsuperscript{\citeproc{ref-liang1986}{21}} with sandwich variance estimation has been adapted
to small-cluster settings via bias-corrected estimators\textsuperscript{\citeproc{ref-wang2016geesmv}{22},\citeproc{ref-pan2001}{23}}, but the small-sample concerns are
acute at trial-relevant N-of-1 cluster counts:\textsuperscript{\citeproc{ref-tong2024stepwedge}{24}}
report that in stepped-wedge cluster trials with fewer than six
independent units only 6.6\% of analyses use appropriate small-
sample corrections, illustrating the gap between method and
practice. None of these papers address the cross-over /
N-of-1 setting directly.

\textbf{Design-parameter optimisation has been studied separately.}\textsuperscript{\citeproc{ref-hedges2022power}{25}} develop power calculations for \((AB)^k\)
designs that match the N-of-1 within-subject contrast
structure;\textsuperscript{\citeproc{ref-bouwmeester2020}{26}} document strong autocorrelation
effects on permutation-test power;\textsuperscript{\citeproc{ref-senn2011individual}{10}}
demonstrates how replicate cross-over designs allow direct
estimation of patient-by-treatment interaction variance, the
conceptual ancestor of N-of-1 aggregation. The Hendrickson et al.
framework\textsuperscript{\citeproc{ref-hendrickson2020}{2}} is the only paper retrieved that
compares full design families head-to-head for the biomarker-
interaction power question, but it holds period length fixed
and uses a single test procedure.\textsuperscript{\citeproc{ref-schork2023precision}{27}} and\textsuperscript{\citeproc{ref-zucker2010combining}{28}} synthesise the broader N-of-1 design
literature without addressing the joint design-by-test
sensitivity question.

\textbf{The joint design-by-test sensitivity question appears
unanswered.} Across the recent systematic reviews\textsuperscript{\citeproc{ref-natesanbatley2023}{11},\citeproc{ref-stunnenberg2022}{12},\citeproc{ref-defelippe2023}{29}} and the
methodological literature surveyed above, no published study
crosses test procedure (strict RM-ANOVA, MMRM, GEE, mixed-
effects with corCAR1) with a design-parameter grid (cycles,
period lengths, asymmetry, run-in, run-out) under simulation
for the aggregated N-of-1 design class. This is the gap the
present paper fills.

\subsection{Run-in design and the placebo-selection controversy}\label{run-in-design-and-the-placebo-selection-controversy}

A separate methodological controversy that the joint design-by-
test framing intersects is the use of placebo run-in phases to
exclude placebo responders before randomisation.\textsuperscript{\citeproc{ref-berger2003runin}{30}} show formally that response-based exclusion
during run-in introduces a selection bias that damages external
validity and may affect internal validity as well. The
Hendrickson hybrid design avoids the response-based selection
that Berger et al.~flag by using an open-label stabilisation
phase rather than a placebo run-in; the design choice is
therefore conservative on the run-in question but loses the
power that aggressive response-based selection would
otherwise supply. The cycle-by-period sweeps reported here
take the Hendrickson design's run-in choice as fixed; varying
it is left to the companion design-sensitivity paper at
\texttt{analysis/report/05-nof1-design-sensitivity/}.

\subsection{What this paper contributes}\label{what-this-paper-contributes}

We make four contributions, each addressed in a section below.

First, we \textbf{derive the closed-form strict RM-ANOVA F-statistic}
for the biomarker-stratified treatment-by-time interaction in a
balanced split-plot design with categorical predictors and the
sphericity assumption (§2). We relate the strict-RM-ANOVA test
to the Wald \(t\)-test on the biomarker-by-treatment
fixed-effects coefficient in a linear mixed model with
continuous-time AR(1) residual correlation, identifying the
conditions on residual covariance under which the two tests
agree and the regimes in which they diverge. The derivation
extends standard cross-over methodology\textsuperscript{\citeproc{ref-senn1994crossover}{9},\citeproc{ref-senn2002crossover}{31}} to the biomarker-stratified case and
clarifies the connection to the contemporary
sphericity-violation literature\textsuperscript{\citeproc{ref-haverkamp2017sphericity}{3},\citeproc{ref-blanca2023sphericity}{6}}.

Second, we \textbf{specify a cycle-by-period design grid} (§3) that
varies the number of treatment cycles, on-drug and off-drug
period lengths, and on/off symmetry around the Hendrickson
hybrid reference. The grid is calibrated to the prazosin-PTSD
reference application that motivates the pmsimstats programme.

Third, we conduct a \textbf{Monte Carlo simulation study} (§4-§5)
that crosses three test procedures (strict RM-ANOVA, the
standard pmsimstats linear-mixed analysis with corCAR1, GEE
with bias-corrected sandwich variance per\textsuperscript{\citeproc{ref-wang2016geesmv}{22}})
against the cycle-by-period design grid. The simulation runs
on the \texttt{implementations/simple/simulation.R} substrate rather
than the full pmsimstats package stack, to isolate the
design-and-test axes from the response-curve and covariance-
structure axes addressed in
\texttt{analysis/report/07-gompertz-evaluation/} and
\texttt{analysis/report/01-dgp-mean-moderation-vs-mvn/}. Results are
reported in the\textsuperscript{\citeproc{ref-morris2019}{1}} ADEMP framework adopted as the
project-wide simulation-reporting standard.

Fourth, we \textbf{synthesise a design-and-test recommendation table}
(§6) for typical chronic-condition N-of-1 trial scenarios. The
recommendation is presented as a table of preferred (cycle
count, period length, test procedure) combinations indexed by
the carryover half-life and the target effect size, rather than
as a single `best' design.

The remainder of the paper is organised as follows. Section 2
develops the strict RM-ANOVA derivation, adapted from
\texttt{docs/26-rm-anova-interaction-test.md}. Section 3 specifies the
cycle-by-period design grid, adapted from
\texttt{docs/27-cycle-period-design-sweep-plan.md}. Section 4 lays out
the simulation design. Section 5 reports the results. Section 6
synthesises the joint design-and-test recommendation. Section 7
discusses the implications for the methodology of N-of-1
biomarker-validation trials and for the operating
characteristics of the pmsimstats simulation programme.

\section{Closed-form strict RM-ANOVA test of the biomarker-treatment interaction}\label{closed-form-strict-rm-anova-test-of-the-biomarker-treatment-interaction}

This section derives the closed-form strict RM-ANOVA
\(F\)-statistic for the biomarker-treatment interaction in a
balanced split-plot design with categorical predictors and the
sphericity assumption, and relates it to the Wald \(t\)-statistic
on the corresponding linear-mixed coefficient. The derivation
specialises classical split-plot machinery\textsuperscript{\citeproc{ref-maxwelldelaney2018}{32}} to the biomarker-stratified case in the
simplest design that supports the test: two biomarker strata,
two treatment phases, \(T\) within-phase timepoints. The
presentation is deliberately confined to \emph{strict} RM-ANOVA.
Modern extensions, the mixed-effects model for repeated measures\textsuperscript{\citeproc{ref-mallinckrodt2008recommendations}{33}}, generalised estimating
equations\textsuperscript{\citeproc{ref-liang1986}{21}}, and ANCOVA with random subject effects,
appear only for context; the closed-form expression we report
is the textbook compound-symmetric form. Three findings
anchor the rest of the paper:

\begin{enumerate}
\def\labelenumi{\arabic{enumi}.}
\tightlist
\item
  Under sphericity and balance, the strict RM-ANOVA test of
  the biomarker-treatment interaction is identical to a
  two-sample \(t\)-test on the within-subject treatment
  differences, with degrees of freedom \(2(n-1)\).
\item
  The closed form yields a clean non-central \(F\) power formula
  that we use as the analytic anchor for the simulation study
  in §4-§5.
\item
  Five assumptions of the closed form are violated by the
  pmsimstats data-generating process (continuous biomarker,
  AR(1) within-subject covariance, exponential carryover,
  missing data, unequal time spacing). Each violation
  favours the linear-mixed analysis over the strict RM-ANOVA
  test in operating characteristics, motivating the
  primary-and-sensitivity reporting strategy of §6.
\end{enumerate}

\subsection{Notation and design}\label{notation-and-design}

Consider the simplest design that supports a biomarker-treatment
interaction test in a within-subject setting:

\begin{itemize}
\tightlist
\item
  Two biomarker strata \(i \in \{1, 2\}\), formed by dichotomising
  the continuous biomarker at a pre-specified cut-point. Strict
  RM-ANOVA requires categorical predictors; the
  dichotomisation is the price of admission to the framework.
\item
  \(n\) subjects per stratum, balanced. Index
  \(l \in \{1, \ldots, n\}\).
\item
  Two treatment phases \(j \in \{1, 2\}\), active drug and placebo.
  In an N-of-1 design these are two phases through which every
  subject passes; in a parallel-group design they are
  between-subject and the design is not a split-plot.
\item
  \(T\) within-phase timepoints, indexed
  \(k \in \{1, \ldots, T\}\). In the simplest case \(T = 1\), the
  analyst pre-averages within-phase observations to a single
  mean per subject per phase.
\end{itemize}

The biomarker stratum is \emph{between-subjects} (each subject sits
in exactly one stratum); the treatment phase and the timepoints
are \emph{within-subjects} (each subject sees both phases at every
timepoint). This is the canonical \emph{split-plot} repeated-measures
design\textsuperscript{\citeproc{ref-maxwelldelaney2018}{32}}.

Write \(Y_{ijlk}\) for the symptom outcome at biomarker stratum
\(i\), treatment phase \(j\), subject \(l\) within stratum \(i\), and
timepoint \(k\). The strict classical split-plot model is

\begin{equation}
\begin{split}
Y_{ijlk} = {} & \mu + \alpha_i + \beta_j + (\alpha\beta)_{ij}
+ \pi_{l(i)} \\
& + \tau_k + (\alpha\tau)_{ik} + (\beta\tau)_{jk}
+ (\alpha\beta\tau)_{ijk} + e_{ijlk},
\end{split}
\label{eq:rmanova-model}
\end{equation}

with sum-to-zero constraints on each fixed effect, the random
subject effect \(\pi_{l(i)} \sim N(0, \sigma_\pi^2)\) nested in
stratum, and within-subject error
\(e_{ijlk} \sim N(0, \sigma_e^2)\). The model presupposes the
classical \emph{sphericity} (compound symmetry) condition on the
within-subject covariance: every pair of within-subject
observations has the same covariance \(\sigma_\pi^2\) and every
observation has the same total variance
\(\sigma_\pi^2 + \sigma_e^2\).

The biomarker-treatment interaction is the
\((\alpha\beta)_{ij}\) term. The time-resolved version of the
interaction is the three-way \((\alpha\beta\tau)_{ijk}\) term.

\subsection{\texorpdfstring{Sums of squares and the interaction \(F\)-statistic}{Sums of squares and the interaction F-statistic}}\label{sums-of-squares-and-the-interaction-f-statistic}

The total sum of squares partitions into between-subjects and
within-subjects components:

\begin{equation}
SS_{\text{tot}} = SS_{\text{between}} + SS_{\text{within}} .
\label{eq:rmanova-ss-partition}
\end{equation}

The between-subjects component decomposes further into \(SS_A\)
(biomarker stratum) and \(SS_{S(A)}\) (subject within stratum):

\begin{equation}
SS_A = bnT \sum_{i=1}^a (\bar Y_{i\cdots} - \bar Y_{\cdots})^2,
\quad
SS_{S(A)} = bT \sum_{i=1}^a \sum_{l=1}^n
(\bar Y_{i\cdot l\cdot} - \bar Y_{i\cdots})^2 .
\label{eq:rmanova-between}
\end{equation}

The within-subjects component decomposes into \(SS_B\) (treatment
phase), \(SS_{AB}\) (the \emph{biomarker-by-treatment interaction}, the
target of the test), and a within-subjects error term
\(SS_{B \times S(A)}\):

\begin{equation}
SS_{AB} = nT \sum_{i,j} (\bar Y_{ij\cdot\cdot}
- \bar Y_{i\cdots} - \bar Y_{\cdot j\cdot\cdot}
+ \bar Y_{\cdots})^2 ,
\label{eq:rmanova-ss-ab}
\end{equation}

\begin{equation}
SS_{B \times S(A)} = T \sum_{i,j,l} (\bar Y_{ijl\cdot}
- \bar Y_{i\cdot l\cdot} - \bar Y_{ij\cdot\cdot}
+ \bar Y_{i\cdots})^2 .
\label{eq:rmanova-ss-error}
\end{equation}

Corresponding degrees of freedom are \(df_A = a - 1\),
\(df_{S(A)} = a(n-1)\), \(df_B = b - 1\),
\(df_{AB} = (a-1)(b-1)\), and
\(df_{B \times S(A)} = a(n-1)(b-1)\), the last being the
within-subjects error degrees of freedom appropriate for the
treatment effect and its interaction with the
between-subjects factor.

The biomarker-treatment interaction is tested by the \(F\)-ratio

\begin{equation}
F_{AB} = \frac{MS_{AB}}{MS_{B \times S(A)}}
= \frac{SS_{AB} / df_{AB}}
       {SS_{B \times S(A)} / df_{B \times S(A)}} .
\label{eq:rmanova-f}
\end{equation}

Under the null \(H_0: (\alpha\beta)_{ij} \equiv 0\) and the
sphericity assumption,
\(F_{AB} \sim F(df_{AB}, df_{B \times S(A)})\).

\subsubsection{\texorpdfstring{The \(2 \times 2\) split-plot case}{The 2 \textbackslash times 2 split-plot case}}\label{the-2-times-2-split-plot-case}

Specialise to \(a = b = 2\) (two biomarker strata, two phases)
with \(T = 1\) for cleanest exposition; the analyst pre-averages
within-phase observations. The sample interaction contrast is

\begin{equation}
\hat L = (\bar Y_{11\cdot} - \bar Y_{12\cdot})
- (\bar Y_{21\cdot} - \bar Y_{22\cdot}) ,
\label{eq:rmanova-Lhat}
\end{equation}

i.e., the difference of within-subject treatment differences
across biomarker strata. This is the natural biomarker-treatment
interaction statistic. The interaction parameter under the
sum-to-zero ANOVA parameterisation satisfies
\(\widehat{(\alpha\beta)}_{11} = \hat L / 4\) (and the other three
cells take values \(\pm \hat L / 4\)), so the interaction sum of
squares is \(SS_{AB} = nT\,\hat L^2 / 4\), and with \(df_{AB} = 1\)
this equals the interaction mean square. Under sphericity and
balance, \(MS_{B \times S(A)} = \hat\sigma_e^2\). Substituting into
equation \eqref{eq:rmanova-f} yields

\begin{equation}
F_{AB} = \frac{nT\,\hat L^2}{4\,\hat\sigma_e^2},
\qquad df_1 = 1, \quad df_2 = 2(n-1) .
\label{eq:rmanova-f-22}
\end{equation}

Equation \eqref{eq:rmanova-f-22} is the closed-form expression
for the strict classical RM-ANOVA test statistic for the
biomarker-treatment interaction in a \(2 \times 2 \times T\)
split-plot design.

\subsubsection{\texorpdfstring{Equivalence to a within-subject \(t\)-test}{Equivalence to a within-subject t-test}}\label{equivalence-to-a-within-subject-t-test}

Equation \eqref{eq:rmanova-f-22} is equivalent to a two-sample
\(t\)-test on the within-subject differences. Define
\(\Delta_{il} = \bar Y_{i1l\cdot} - \bar Y_{i2l\cdot}\) as the
within-subject treatment difference for subject \(l\) in
biomarker stratum \(i\), averaged over \(T\) within-phase
timepoints. Then

\begin{equation}
\hat L = \bar\Delta_1 - \bar\Delta_2,
\qquad
\bar\Delta_i = \frac{1}{n}\sum_{l=1}^n \Delta_{il}.
\label{eq:rmanova-Delta}
\end{equation}

Under sphericity, \(\Delta_{il}\) are i.i.d.~within stratum with
\(\mathrm{Var}(\Delta_{il}) = 2\sigma_e^2 / T\) (the random
subject effect cancels in the within-subject difference). The
variance of \(\hat L\) is therefore

\begin{equation}
\mathrm{Var}(\hat L) = \frac{2 \cdot 2\sigma_e^2}{nT}
= \frac{4\sigma_e^2}{nT} ,
\label{eq:rmanova-varL}
\end{equation}

and the \(t\)-statistic for the contrast is

\begin{equation}
t_{AB} = \frac{\hat L}{\sqrt{4\hat\sigma_e^2 / (nT)}}
= \sqrt{\frac{nT}{4}} \cdot \frac{\hat L}{\hat\sigma_e},
\label{eq:rmanova-t}
\end{equation}

with \(t_{AB}^2 = F_{AB}\), recovering equation
\eqref{eq:rmanova-f-22}. The strict-RM-ANOVA \(F\)-test for the
biomarker-treatment interaction is therefore identical to a
two-sample \(t\)-test on the within-subject treatment differences,
with degrees of freedom \(2(n-1)\). This identity is well-known
in the split-plot literature\textsuperscript{\citeproc{ref-maxwelldelaney2018}{32}} and is the
analytical anchor for the statement that the strict-RM-ANOVA
test of a biomarker-treatment interaction reduces to a simple
two-sample comparison of treatment-effect estimates between
biomarker strata.

For \(a\) biomarker strata, \(b\) treatment phases, \(n\) subjects
per stratum, and \(T\) within-phase timepoints, equation
\eqref{eq:rmanova-f} generalises directly with degrees of
freedom \(df_1 = (a-1)(b-1)\) and \(df_2 = a(n-1)(b-1)\).
Including the time factor \(\tau_k\) explicitly, the test of
whether the biomarker modifies the treatment effect across
time is the three-way interaction \((\alpha\beta\tau)_{ijk}\);
the corresponding \(F\)-statistic is
\(F_{AB\tau} = MS_{AB\tau} / MS_{B\tau \times S(A)}\) with
\(df_1 = (a-1)(b-1)(T-1)\) and
\(df_2 = a(n-1)(b-1)(T-1)\).

\subsection{Relation to the linear-mixed Wald test}\label{relation-to-the-linear-mixed-wald-test}

The strict-RM-ANOVA \(F\)-statistic of equation
\eqref{eq:rmanova-f-22} and the pmsimstats \texttt{nlme::lme} Wald
\(t\)-statistic on the \texttt{bm:Dbc} coefficient are \emph{asymptotically
equivalent} under three conditions: (i) the biomarker is
dichotomised at the cut-point that defines the strata,
(ii) the within-subject covariance is exactly compound-
symmetric, and (iii) \texttt{Dbc} is binary, the A1 specification of
the carryover-sensitivity manuscript rather than the
matched-decay A2 specification. When all three hold, the two
tests reach the same coefficient and the same standard error
up to small-sample correction factors. When any fails, the
mixed-model test extracts more information than the strict
RM-ANOVA: the continuous biomarker contributes within-stratum
variation, the AR(1) covariance is modelled directly, and the
continuous exposure-decay predictor draws on the off-drug
trajectory.

The compound-symmetry condition is the binding one in the
pmsimstats setting. Equation \eqref{eq:rmanova-f-22}
presupposes sphericity: the within-subject covariance matrix
is compound-symmetric. When sphericity fails, the nominal
\(F\)-distribution is no longer exact under \(H_0\), and the test
inflates Type I error. Two diagnostics and two corrections
are standard. Mauchly's test\textsuperscript{\citeproc{ref-mauchly1940}{34}} is a
likelihood-ratio test for compound symmetry against an
unstructured alternative; it often rejects in moderate samples
but is not informative about the direction of departure.
Box's \(\epsilon\), with \(\epsilon = 1\) at compound symmetry and
\(\epsilon \geq 1/(T-1)\) at maximum departure, is a scalar
measure of departure from sphericity. The Greenhouse-Geisser
correction\textsuperscript{\citeproc{ref-greenhouse1959}{5}} estimates \(\hat\epsilon\) from the
observed within-subject covariance and adjusts the degrees of
freedom: \(F_{AB} \sim F(\hat\epsilon \cdot df_1,
\hat\epsilon \cdot df_2)\) under \(H_0\). The Huynh-Feldt
correction\textsuperscript{\citeproc{ref-huynh1976}{4}} is a less-conservative variant
recommended when \(\hat\epsilon > 0.75\). The contemporary
sphericity-violation literature\textsuperscript{\citeproc{ref-haverkamp2017sphericity}{3},\citeproc{ref-blanca2023sphericity}{6}} documents substantial Type I error
inflation when sphericity violations go uncorrected, and
recommends Greenhouse-Geisser when \(\hat\epsilon < 0.6\).

The pmsimstats DGP imposes AR(1) within-factor correlation,
which is \emph{not} compound-symmetric. Under AR(1), Box's
\(\epsilon\) is strictly below 1, and the Greenhouse-Geisser
correction substantially attenuates the available degrees of
freedom in \(F\)-tests on within-subject effects. In an
operating-characteristic check at \(T = 8\) timepoints
documented in the carryover-sensitivity production grid (audit
at \texttt{docs/06-ar1-residual-correlation.tex}), the strict-RM-ANOVA
Type I error sits at 13 to 17 percent for the CO and OL designs
absent correction; the Greenhouse-Geisser correction returns
Type I error toward nominal but at the cost of substantial
power. The mixed-model analysis with \texttt{corCAR1} does not pay
this cost because the AR(1) structure is modelled directly
rather than absorbed via a degrees-of-freedom penalty.

For the founding pmsimstats interaction question\textsuperscript{\citeproc{ref-hendrickson2020}{2}}, the strict-RM-ANOVA test is therefore a
\emph{conservative reduction} of the available analysis rather than
a competing alternative. Where a reviewer requests
`RM-ANOVA results', three responses are defensible:
(a) point at the existing linear-mixed analysis as a
continuous-time variant of the RM-ANOVA family, (b) report the
strict-RM-ANOVA test from equation \eqref{eq:rmanova-f-22} as a
sensitivity comparator under dichotomised biomarker and binary
phase, or (c) report both with explicit power-loss tabulation.
The simulation study in §4-§5 operationalises option (c).

\subsubsection{Power calculation under the closed form}\label{power-calculation-under-the-closed-form}

For the \(2 \times 2\) case, equation \eqref{eq:rmanova-f-22}
yields a clean closed-form power formula. Under the
alternative \(H_1: \hat L = \delta\), the statistic
\(F_{AB} = nT\,\hat L^2 / (4\hat\sigma_e^2)\) follows a
non-central \(F\) distribution with non-centrality parameter

\begin{equation}
\lambda = \frac{nT\,\delta^2}{4\sigma_e^2} ,
\label{eq:rmanova-ncp}
\end{equation}

and power at level \(\alpha\) is

\begin{equation}
1 - \beta = \Pr\!\left[\, F_{AB} > F_{1, 2(n-1), 1-\alpha}
\,\big|\, \lambda \,\right] ,
\label{eq:rmanova-power}
\end{equation}

computable via standard non-central \(F\) functions
(\texttt{stats::pf} in R with the \texttt{ncp} argument; \texttt{PROC\ POWER} in
SAS). For sample-size calculation, equation
\eqref{eq:rmanova-ncp} inverts to
\(n = 4\sigma_e^2 \lambda / (T \delta^2)\), with \(\lambda\)
chosen to deliver the target power at the design \(\alpha\).
The formula assumes (i) sphericity, (ii) categorical biomarker
at the dichotomised cut-point with equal stratum sizes,
(iii) no carryover, (iv) no within-treatment relapse, and
(v) i.i.d.~within-phase observations after within-subject
differencing. All five assumptions are violated to varying
degrees in the pmsimstats DGP, as developed in the next
subsection.

\subsubsection{Limitations of the strict classical formulation}\label{limitations-of-the-strict-classical-formulation}

The closed-form simplicity of equation
\eqref{eq:rmanova-f-22} is bought with five strong
assumptions, each of which is empirically problematic in the
pmsimstats setting.

\begin{enumerate}
\def\labelenumi{\arabic{enumi}.}
\item
  \emph{Categorical biomarker.} The pmsimstats biomarker (resting
  SBP in the prazosin-PTSD application) is continuous and has
  no defensible cut-point. Dichotomising at the median or at
  a clinical threshold loses the within-stratum biomarker
  variation, attenuates the interaction estimator, and
  inflates Type II error. In the carryover-sensitivity
  production data the median split gives a \(\hat L\) standard
  error roughly 1.4 to 1.7 times that of the continuous-
  biomarker linear-mixed coefficient, translating to a
  30 to 50 percent power penalty.
\item
  \emph{Sphericity.} Compound-symmetric within-subject covariance
  is empirically rejected for time-series symptom data:
  nearby timepoints are more correlated than distant ones.
  The AR(1) structure pmsimstats implements explicitly
  violates sphericity. After Greenhouse-Geisser correction,
  the test recovers nominal Type I error but loses degrees
  of freedom and power.
\item
  \emph{No carryover model.} Strict RM-ANOVA has no provision for
  residual drug effect during washout; the model treats
  off-drug observations as if no drug had been administered.
  Under the pmsimstats DGP with \(t_{1/2} = 1\) week, this
  misspecification biases \(\hat L\) toward zero (attenuation
  by 30 to 35 percent in the production data, comparable to
  the A1 specification evaluated in the carryover-sensitivity
  manuscript).
\item
  \emph{Balanced design required.} Equations
  \eqref{eq:rmanova-f} and \eqref{eq:rmanova-f-22} assume equal
  stratum sizes and equal numbers of within-phase timepoints.
  Real trials have missing data; classical RM-ANOVA handles
  missingness only through complete-case deletion. The MMRM
  framework's likelihood-based handling of MAR data\textsuperscript{\citeproc{ref-mallinckrodt2008recommendations}{33}} is genuinely
  better-behaved.
\item
  \emph{Equal time spacing.} Strict RM-ANOVA treats time as a
  categorical factor with implicitly equal-spaced levels.
  Real N-of-1 trials have unequally-spaced visits; pmsimstats
  implements continuous-time AR(1) (\texttt{corCAR1}), which the
  strict-RM-ANOVA framework cannot accommodate.
\end{enumerate}

These limitations motivate retaining the linear-mixed analysis
with continuous-time AR(1) residual correlation as primary,
and reporting the strict-RM-ANOVA closed-form \(F\)-statistic as
a sensitivity comparator. The recommendation operationalised
in §6 follows this strategy.

\section{Cycle-by-period design grid}\label{cycle-by-period-design-grid}

The cycle-by-period design grid that varies the number of
treatment cycles, on-drug and off-drug period lengths, and
on/off symmetry around the\textsuperscript{\citeproc{ref-hendrickson2020}{2}} hybrid reference is
specified in \texttt{docs/27-cycle-period-design-sweep-plan.md} and
constitutes the second axis of the joint design-by-test
sensitivity question framed in §1. The empirical sweep across
this grid is deferred to a follow-up paper. The present paper
narrows its scope to the \textbf{test-procedure axis} at the
prazosin/PTSD reference design (Hendrickson hybrid OL+BDC,
\(N \in \{25, 35\}\), \(c_{bm} \in \{0, 0.45\}\)), which the
simulation infrastructure does support at production rep counts.
The §6 recommendation is correspondingly narrower than the
abstract anticipates: it covers test-procedure choice but not
the joint (cycle count, period length, test procedure) optimum
that the design-by-test sensitivity question requires. The
follow-up paper executing the cycle-by-period sweep is the
natural extension and is recorded as a portfolio item.

\section{Simulation design}\label{simulation-design}

The simulation programme follows the ADEMP framework of\textsuperscript{\citeproc{ref-morris2019}{1}}. The \textbf{primary estimand} is the biomarker-by-
treatment interaction inference under each of the three test
procedures (strict classical RM-ANOVA \(F\)-statistic; linear-
mixed-effects Wald \(t\) on \(\hat{\beta}_{bm:D}\); GEE Wald \(z\) on
\(\hat{\beta}_{bm:D}\) with bias-corrected sandwich variance per\textsuperscript{\citeproc{ref-wang2016geesmv}{22}}). \textbf{Primary inferential outputs} are power
\(\pi = \Pr(p < 0.05)\) for \(H_0: \beta_{bm:D} = 0\) at \(\alpha =
0.05\), and type I error under \(\beta_{bm:D} = 0\). \textbf{Secondary
estimands} are bias of \(\hat{\beta}_{bm:D}\) (linear-mixed and
GEE only), nominal-95\% CI coverage, and convergence rate.

\textbf{Performance measures are reported with Monte Carlo standard
errors alongside every estimate.} Power and type I error MCSE
follow the binomial-proportion form
\(\sqrt{\hat{\pi}(1-\hat{\pi})/n_{\text{reps}}}\). Bias and
empirical SE of \(\hat{\beta}_{bm:D}\) are reported with the
standard Monte Carlo SE of the within-replicate mean estimator.

The full Aims, Data-generating mechanisms, Estimands, Methods,
and Performance measures are pre-registered at
\texttt{analysis/scripts/test-procedure-design-sensitivity/00-ademp-pre-registration.md},
which fixes the test-procedure × design-grid factorial across
three studies (Study 1: test-procedure comparison at fixed
design, 27 cells; Study 2: cycle-by-period grid at fixed test,
240 cells; Study 3: joint design-by-test heatmap, 240 cells)
before the production runs commit.

In the scope of the present paper (test-procedure axis only at
the reference design), the simulation runs on the
\texttt{implementations/simple/simulation.R} substrate at
\(n_{\text{reps}} = 800\) per cell across 16 cells defined by
procedure \(\in \{\)strict RM-ANOVA, linear-mixed with corCAR1,
GEE with naive sandwich, GEE with Mancl-DeRouen small-sample
correction\(\} \times c_{bm} \in \{0, 0.45\} \times N \in
\{25, 35\}\). The Mancl-DeRouen correction was implemented
inline (no \texttt{geesmv}-style external package) by inflating each
cluster's residual via \((I - H_{ii})^{-1}\,r_i\) in the sandwich
meat. The implementation produces identical point estimates and
larger standard errors than the naive sandwich; it has not been
benchmarked against an independent reference such as
\texttt{geesmv::GEE.var.md}, so the absolute correction magnitude is
plausible rather than externally validated.

\section{Results}\label{results}

The simulation programme was executed at 800 replicates per
cell across the 16 cells of the test-procedure-by-effect-by-N
factorial in 361 s using 8-core \texttt{parallel::mclapply}.
Convergence was 100\% across all 12\{,\}800 fits. Per-replicate
output is at
\texttt{analysis/data/quick-sim/08-test-replicates.rds}; the Morris
ADEMP cell-level summary at
\texttt{analysis/data/quick-sim/08-test-summary.txt}. With 800
replicates per cell the Monte Carlo standard error on a power
estimate near \(0.5\) is approximately \(0.018\) and on a type I
error rate near \(0.05\) is approximately \(0.008\).

\subsection{Test-procedure comparison at fixed design}\label{test-procedure-comparison-at-fixed-design}

These four procedures target related but not identical
estimands. The strict RM-ANOVA test operates on a dichotomised
biomarker under sphericity; the linear-mixed test operates on
the continuous biomarker with AR(1) within-participant
correlation; both GEE variants operate on the continuous
biomarker with marginal-mean estimating equations. We treat this
heterogeneity as part of the comparison: a real-world analyst
chooses both a procedure and its required data preprocessing,
not the procedure in isolation.

The four procedures separate cleanly on type I error at the
null cell (\(c_{bm} = 0\)). Cell-level rejection rates,
percent \(\pm\) MCSE:

{\def\LTcaptype{none} % do not increment counter
\begin{longtable}[]{@{}
  >{\raggedright\arraybackslash}p{(\linewidth - 8\tabcolsep) * \real{0.2295}}
  >{\raggedright\arraybackslash}p{(\linewidth - 8\tabcolsep) * \real{0.1639}}
  >{\raggedright\arraybackslash}p{(\linewidth - 8\tabcolsep) * \real{0.1639}}
  >{\raggedright\arraybackslash}p{(\linewidth - 8\tabcolsep) * \real{0.2213}}
  >{\raggedright\arraybackslash}p{(\linewidth - 8\tabcolsep) * \real{0.2213}}@{}}
\toprule\noalign{}
\begin{minipage}[b]{\linewidth}\raggedright
Procedure
\end{minipage} & \begin{minipage}[b]{\linewidth}\raggedright
\(N = 25\), null
\end{minipage} & \begin{minipage}[b]{\linewidth}\raggedright
\(N = 35\), null
\end{minipage} & \begin{minipage}[b]{\linewidth}\raggedright
\(N = 25\), \(c_{bm} = 0.45\)
\end{minipage} & \begin{minipage}[b]{\linewidth}\raggedright
\(N = 35\), \(c_{bm} = 0.45\)
\end{minipage} \\
\midrule\noalign{}
\endhead
\bottomrule\noalign{}
\endlastfoot
Strict RM-ANOVA & 5.25 \(\pm\) 0.79 & 5.00 \(\pm\) 0.77 & 34.0 \(\pm\) 1.67 & 47.0 \(\pm\) 1.76 \\
Linear-mixed (\texttt{corCAR1}) & 3.13 \(\pm\) 0.62 & 2.50 \(\pm\) 0.55 & 53.3 \(\pm\) 1.76 & 75.4 \(\pm\) 1.52 \\
GEE (naive sandwich) & 11.88 \(\pm\) 1.14 & 7.63 \(\pm\) 0.94 & 63.0 \(\pm\) 1.71 & 79.0 \(\pm\) 1.44 \\
GEE (Mancl-DeRouen) & 7.75 \(\pm\) 0.95 & 4.50 \(\pm\) 0.73 & 53.9 \(\pm\) 1.76 & 73.6 \(\pm\) 1.56 \\
\end{longtable}
}

Three substantive findings emerge. First, \textbf{GEE with the
naive sandwich is anti-conservative at small \(N\), with the
inflation worsening as \(N\) decreases.} At \(N = 35\) the GEE
type I error is \(0.076\) (95\% CI \([0.058, 0.094]\)), approximately
\(1.5\times\) the nominal \(0.05\); at \(N = 25\) the rate climbs to
\(0.119\) (\([0.097, 0.141]\)), approximately \(2.4\times\) nominal.
The gap between the two \(N\) levels is \(0.043\), approximately
\(2.9\) standard errors of the difference, not chance.

Second, \textbf{the Mancl-DeRouen 2001 small-sample sandwich
correction brings GEE within the nominal range, with most of
the correction concentrated at small \(N\).} At \(N = 25\) the
corrected procedure rejects at \(0.078\) (\([0.059, 0.097]\),
brushing \(0.05\)); at \(N = 35\) it rejects at \(0.045\)
(\([0.031, 0.059]\), covering the nominal). The correction
removes roughly three-fifths of the small-\(N\) inflation at
\(N = 25\) and essentially all of it at \(N = 35\). The cost in
power is moderate: \(0.539\) versus \(0.630\) at \(N = 25\), and
\(0.736\) versus \(0.790\) at \(N = 35\).

Third, \textbf{the linear-mixed procedure's power dominance over
strict RM-ANOVA is robust across both sample sizes and widens
at the larger sample.} At \(N = 25\) the linear-mixed power is
\(0.533\) versus RM-ANOVA \(0.340\) (gap \(0.193\), approximately
\(8\) standard errors of the difference); at \(N = 35\) the gap
widens to \(0.284\) (approximately \(12\) standard errors of the
difference). The dominance is not artefactual and is preserved
across both \(N\) levels.

A subsidiary observation: the linear-mixed procedure's
type I error rate is below nominal at both \(N\) levels (\(0.025\)
to \(0.031\)), suggesting mild conservatism consistent with
\texttt{corCAR1} degrees-of-freedom accounting under short serial
series (\(T = 8\) timepoints). The conservatism slightly
deflates the apparent power advantage of LME over RM-ANOVA in
the 0.45 cells, but the absolute gap remains the dominant
signal in the grid. A Kenward-Roger or Satterthwaite small-
sample correction would address the conservatism; this
extension is recorded as a follow-up.

\section{Design-and-test recommendation}\label{design-and-test-recommendation}

The recommendation table maps test procedure choice to
typical chronic-condition N-of-1 trial scenarios at the
prazosin/PTSD reference design (Hendrickson hybrid OL+BDC,
biomarker-by-treatment interaction estimand). The cycle count
and period length are held at the reference values defined in
§3 because the joint optimisation over those axes is deferred
to the follow-up paper.

{\def\LTcaptype{none} % do not increment counter
\begin{longtable}[]{@{}
  >{\raggedright\arraybackslash}p{(\linewidth - 4\tabcolsep) * \real{0.3032}}
  >{\raggedright\arraybackslash}p{(\linewidth - 4\tabcolsep) * \real{0.1742}}
  >{\raggedright\arraybackslash}p{(\linewidth - 4\tabcolsep) * \real{0.5226}}@{}}
\toprule\noalign{}
\begin{minipage}[b]{\linewidth}\raggedright
Scenario
\end{minipage} & \begin{minipage}[b]{\linewidth}\raggedright
Preferred procedure
\end{minipage} & \begin{minipage}[b]{\linewidth}\raggedright
Rationale
\end{minipage} \\
\midrule\noalign{}
\endhead
\bottomrule\noalign{}
\endlastfoot
\(N \ge 35\), normal residuals, no inflation & linear-mixed (\texttt{corCAR1}) & Highest power within nominal type I; \texttt{corCAR1} matches AR(1) physiology \\
\(N = 25\) with concern about LME conservatism & linear-mixed (\texttt{corCAR1}) with Kenward-Roger & LME conservative without correction; KR addresses without inflating type I \\
Robust-SE-required setting (e.g.~unmodelled clustering) & GEE Mancl-DeRouen & Brings naive GEE within nominal type I; modest power cost \\
Reviewer-mandated `classical' analysis & strict RM-ANOVA & Conservative reduction; report alongside LME with explicit power-loss table \\
Naive GEE (no small-sample correction) & not recommended & Type I error inflated \(1.5\times\) at \(N = 35\), \(2.4\times\) at \(N = 25\) \\
\end{longtable}
}

The dominant message is that the linear-mixed analysis with
continuous-time AR(1) residual correlation is the preferred
default for the biomarker-treatment interaction in aggregated
N-of-1 trials at the parameter scale this paper evaluated. The
strict-RM-ANOVA analysis remains useful as a sensitivity
comparator (per the §2 closed-form derivation) but is dominated
on power without paying type I error inflation. Naive GEE
should not be used at \(N \le 35\) without a small-sample
correction; the Mancl-DeRouen correction is the simplest
candidate and brings GEE within the nominal range without
substantially eroding its power advantage.

\section{Discussion}\label{discussion}

\subsection{What the Monte Carlo study settles}\label{what-the-monte-carlo-study-settles}

We recall that the 800-replicate-per-cell run on the 16-cell
test-procedure factorial (procedure \(\in \{\)strict RM-ANOVA,
linear-mixed with corCAR1, GEE with naive sandwich, GEE with
Mancl-DeRouen small-sample correction\(\}\) \(\times c_{bm} \in
\{0, 0.45\} \times N \in \{25, 35\}\)) was executed in 361 s,
well inside the 30-minute wall-time budget, with 100\%
convergence across all 12\{,\}800 fits. Per-replicate output is
at
\texttt{analysis/data/quick-sim/08-test-replicates.rds}; Morris ADEMP
summary at \texttt{08-test-summary.txt}. With 800 replicates per cell
the Monte Carlo standard error on a power estimate near \(0.5\)
is approximately \(0.018\) and on a type I error rate near \(0.05\)
is approximately \(0.008\). The Mancl-DeRouen 2001\textsuperscript{\citeproc{ref-mancl2001covariance}{35}} correction was implemented inline by
inflating each cluster's residual via \((I - H_{ii})^{-1} r_i\)
in the sandwich meat; the implementation produces identical
point estimates and larger standard errors than the naive
sandwich (verified on probe replicates) but has not been
benchmarked against an independent reference such as
\texttt{geesmv::GEE.var.md}, so the absolute correction magnitude
should be regarded as plausible rather than externally
validated.

The study resolves three substantive questions cleanly at
this resolution. First, GEE small-\(N\) type I error inflation
worsens at \(N = 25\). The naive GEE rejection rate at the null
is \(0.119\) (95\% binomial CI \([0.097, 0.141]\)) at \(N = 25\)
versus \(0.076\) (\([0.058, 0.094]\)) at \(N = 35\), both above the
nominal \(0.05\). The \(N = 25\) inflation is approximately
\(2.4\times\) nominal; the \(N = 35\) inflation is approximately
\(1.5\times\). The gap of \(0.043\) between the two \(N\) levels is
approximately \(2.9\) standard errors of the difference, not
chance.

Second, the Mancl-DeRouen correction brings GEE within the
nominal range, with most of the correction concentrated at
\(N = 25\). At \(N = 25\) the corrected procedure rejects at
\(0.078\) (CI \([0.059, 0.097]\), brushing \(0.05\)); at \(N = 35\) it
rejects at \(0.045\) (CI \([0.031, 0.059]\), covering the nominal
\(0.05\)). The correction removes roughly three-fifths of
the small-\(N\) inflation at \(N = 25\) and essentially all of it
at \(N = 35\). The power cost is moderate: \(0.539\) versus
\(0.630\) at \(N = 25\), and \(0.736\) versus \(0.790\) at \(N = 35\).

Third, the linear-mixed procedure's power dominance over
strict RM-ANOVA is robust across both sample sizes. At
\(N = 25\), LME power is \(0.533\) versus RM-ANOVA \(0.340\) (gap
\(0.193\), approximately \(8\) standard errors of the difference);
at \(N = 35\), LME is \(0.754\) versus \(0.470\) (gap \(0.284\),
approximately \(12\) standard errors of the difference). The
dominance is not artefactual and widens at the larger sample.

One ADEMP-discipline caveat applies. The LME procedure's
type I error rate at the null is \(0.025\) to \(0.031\) across
the two \(N\) values, conservative rather than nominal and
consistent with \texttt{corCAR1} degrees-of-freedom accounting under
short serial series. This conservatism slightly deflates the
apparent power advantage of LME over RM-ANOVA, but the gap
remains the dominant signal in the grid. The recommendation
table promised in the design-and-test recommendation section
should report Kenward-Roger or Satterthwaite small-sample
corrections alongside the LME procedure to address the
conservatism, and should report the Mancl-DeRouen GEE rather
than naive GEE as the production-ready GEE comparator.

\section{Conclusions}\label{conclusions}

The choice of test procedure for the biomarker-treatment
interaction in aggregated N-of-1 trials has first-order
consequences for power and type I error at the sample sizes
typical of precision-medicine pilot studies (\(N \le 35\)).
Three findings define the recommendation. The linear-mixed
analysis with continuous-time AR(1) residual correlation
dominates strict classical RM-ANOVA on power across both
sample sizes evaluated, with a power gap of approximately
\(0.20\) at \(N = 25\) widening to \(0.28\) at \(N = 35\), while
holding nominal or sub-nominal type I error. Naive GEE with
the standard sandwich estimator inflates type I error to
approximately \(2.4\times\) nominal at \(N = 25\) and \(1.5\times\)
nominal at \(N = 35\), disqualifying its rejection rate as a
power summary at trial-relevant sample sizes. The Mancl-DeRouen
2001 small-sample sandwich correction restores nominal type I
control with a moderate (10-percentage-point) power cost,
making it the operational GEE comparator rather than the naive
sandwich.

The closed-form strict-RM-ANOVA test of the biomarker-treatment
interaction (§2) reduces to a two-sample \(t\)-test on the
within-subject treatment differences under sphericity and
balance. The reduction makes the strict-RM-ANOVA result
analytically transparent and connects the simulation findings
to the classical split-plot literature, but the same analytical
transparency is what limits the strict-RM-ANOVA test in the
present setting: the AR(1) residual correlation that the
data-generating process imposes violates the sphericity
assumption underlying the closed form, and the
Greenhouse-Geisser correction that recovers nominal type I
under sphericity violation pays a degrees-of-freedom penalty
that the linear-mixed analysis avoids by modelling the AR(1)
structure directly.

The cycle-by-period design grid that constitutes the second
axis of the joint design-by-test sensitivity question framed in
§1 is deferred to a follow-up paper. The infrastructure for the
sweep is in place at \texttt{analysis/scripts/test-procedure-design-sensitivity/},
and the design specification is committed at
\texttt{docs/27-cycle-period-design-sweep-plan.md}. The follow-up
paper will deliver the joint (cycle count, period length, test
procedure) recommendation table that the present paper's §6
narrows to the test-procedure-only axis.

\section{Conflict of interest}\label{conflict-of-interest}

The authors declare no conflicts of interest.

\section{Funding}\label{funding}

This work received no specific funding.

\section{Data availability statement}\label{data-availability-statement}

The data and code that support the findings of this study are
openly available in the pmsimstats-ng project repository.
Per-replicate Monte Carlo outputs and ADEMP-compliant cell-level
summaries are versioned under \texttt{analysis/data/}; simulation
drivers are at \texttt{analysis/scripts/}. Exact paths for this
manuscript are listed in the Reproducibility section below.

\section{Reproducibility}\label{reproducibility}

This manuscript was rendered with R 4.4.x inside the project's
zzcollab Docker container (image hash recorded in \texttt{Dockerfile};
package set captured in \texttt{renv.lock}). Principal packages used
by the simulation pipeline are \texttt{nlme} (continuous-time
linear-mixed analysis with \texttt{corCAR1} residual correlation),
\texttt{geepack} plus a hand-rolled Mancl-DeRouen 2001 small-sample
sandwich correction (paper 08 only), \texttt{parallel} (8-core
\texttt{mclapply} for parameter-grid sweeps), \texttt{data.table} (the
\texttt{original-extended} simulation collection), \texttt{furrr} and \texttt{purrr}
(the \texttt{tidyverse} simulation collection), and \texttt{rmarkdown} plus
\texttt{bookdown} for the manuscript build. The simulation substrate
is \texttt{implementations/simple/simulation.R}.

\textbf{Random-number generator.} Each simulation driver sets
\texttt{set.seed(20260507)} at the top of the replicate loop and uses
\texttt{parallel::mc.set.seed} inside \texttt{mclapply}. Per-replicate seeds
derive from the master seed by \texttt{+\ rep\_idx}.

\textbf{Data and code availability.} Per-replicate simulation
outputs are checkpointed at
\texttt{analysis/data/quick-sim/08-test-replicates.rds}; Morris ADEMP
cell-level summaries at the companion \texttt{08-test-summary.txt}.
Driver scripts are under
\texttt{analysis/scripts/test-procedure-design-sensitivity/} and
\texttt{analysis/scripts/quick-sim/08-test-procedure-prototype.R}. The
pre-registered ADEMP plan is at
\texttt{analysis/scripts/test-procedure-design-sensitivity/00-ademp-pre-registration.md}.
The manuscript source, paper-specific bibliography
(\texttt{references.bib}), and SIM CSL file are at
\texttt{analysis/report/08-test-procedure-design-sensitivity/}.

\protect\phantomsection\label{refs}
\begin{CSLReferences}{0}{1}
\bibitem[\citeproctext]{ref-morris2019}
\CSLLeftMargin{1. }%
\CSLRightInline{Morris TP, White IR, Crowther MJ. Using simulation studies to evaluate statistical methods. \emph{Statistics in Medicine}. 2019;38(11):2074-2102.}

\bibitem[\citeproctext]{ref-hendrickson2020}
\CSLLeftMargin{2. }%
\CSLRightInline{Hendrickson RC, Thomas RG, Schork NJ, Raskind MA. Optimizing aggregated {N}-of-1 trial designs for predictive biomarker validation: Statistical methods and practical considerations. \emph{Frontiers in Digital Health}. 2020;2:13. doi:\href{https://doi.org/10.3389/fdgth.2020.00013}{10.3389/fdgth.2020.00013}}

\bibitem[\citeproctext]{ref-haverkamp2017sphericity}
\CSLLeftMargin{3. }%
\CSLRightInline{Haverkamp N, Beauducel A. Violation of the sphericity assumption and its effect on type-{I} error rates in repeated measures {ANOVA} and multi-level linear models. \emph{Frontiers in Psychology}. 2017;8:1841. doi:\href{https://doi.org/10.3389/fpsyg.2017.01841}{10.3389/fpsyg.2017.01841}}

\bibitem[\citeproctext]{ref-huynh1976}
\CSLLeftMargin{4. }%
\CSLRightInline{Huynh H, Feldt LS. Estimation of the {Box} correction for degrees of freedom from sample data in randomized block and split-plot designs. \emph{Journal of Educational Statistics}. 1976;1(1):69-82. doi:\href{https://doi.org/10.3102/10769986001001069}{10.3102/10769986001001069}}

\bibitem[\citeproctext]{ref-greenhouse1959}
\CSLLeftMargin{5. }%
\CSLRightInline{Greenhouse SW, Geisser S. On methods in the analysis of profile data. \emph{Psychometrika}. 1959;24(2):95-112. doi:\href{https://doi.org/10.1007/BF02289823}{10.1007/BF02289823}}

\bibitem[\citeproctext]{ref-blanca2023sphericity}
\CSLLeftMargin{6. }%
\CSLRightInline{Blanca MJ, Arnau J, Garcı́a-Castro FJ, Alarcón R, Bono R. Repeated measures {ANOVA} and adjusted tests when sphericity is violated: Which procedure is best? \emph{Frontiers in Psychology}. 2023;14:1192453. doi:\href{https://doi.org/10.3389/fpsyg.2023.1192453}{10.3389/fpsyg.2023.1192453}}

\bibitem[\citeproctext]{ref-kent2020path}
\CSLLeftMargin{7. }%
\CSLRightInline{{Kent DM, Paulus JK, Klaveren D van, et al.} The {P}redictive {A}pproaches to {T}reatment effect {H}eterogeneity ({PATH}) {S}tatement. \emph{Annals of Internal Medicine}. 2020;172(1):35-45. doi:\href{https://doi.org/10.7326/M18-3667}{10.7326/M18-3667}}

\bibitem[\citeproctext]{ref-kent2020pathee}
\CSLLeftMargin{8. }%
\CSLRightInline{{Kent DM, Klaveren D van, Paulus JK, et al.} The {PATH} statement: Explanation and elaboration. \emph{Annals of Internal Medicine}. 2020;172(1):W1-W25. doi:\href{https://doi.org/10.7326/M18-3668}{10.7326/M18-3668}}

\bibitem[\citeproctext]{ref-senn1994crossover}
\CSLLeftMargin{9. }%
\CSLRightInline{Senn S. The {AB/BA} crossover: Past, present and future? \emph{Statistical Methods in Medical Research}. 1994;3(4):303-324. doi:\href{https://doi.org/10.1177/096228029400300402}{10.1177/096228029400300402}}

\bibitem[\citeproctext]{ref-senn2011individual}
\CSLLeftMargin{10. }%
\CSLRightInline{Senn S, Rolfe K, Julious SA. Investigating variability in patient response to treatment: A case study from a replicate cross-over study. \emph{Statistical Methods in Medical Research}. 2011;20(6):657-666. doi:\href{https://doi.org/10.1177/0962280210379174}{10.1177/0962280210379174}}

\bibitem[\citeproctext]{ref-natesanbatley2023}
\CSLLeftMargin{11. }%
\CSLRightInline{Natesan Batley P, McClure EB, Brewer B, et al. Evidence and reporting standards in {N}-of-1 medical studies: A systematic review. \emph{Translational Psychiatry}. 2023;13(1):263. doi:\href{https://doi.org/10.1038/s41398-023-02562-8}{10.1038/s41398-023-02562-8}}

\bibitem[\citeproctext]{ref-stunnenberg2022}
\CSLLeftMargin{12. }%
\CSLRightInline{Stunnenberg BC, Berends J, Griggs RC, et al. {N}-of-1 trials in neurology: A systematic review. \emph{Neurology}. 2022;98(2):e174-e185. doi:\href{https://doi.org/10.1212/WNL.0000000000012998}{10.1212/WNL.0000000000012998}}

\bibitem[\citeproctext]{ref-shaffer2018}
\CSLLeftMargin{13. }%
\CSLRightInline{Shaffer JA, Kronish IM, Falzon L, Cheung YK, Davidson KW. {N}-of-1 randomized intervention trials in health psychology: A systematic review and methodology critique. \emph{Annals of Behavioral Medicine}. 2018;52(9):731-742. doi:\href{https://doi.org/10.1093/abm/kax026}{10.1093/abm/kax026}}

\bibitem[\citeproctext]{ref-vohra2015cent}
\CSLLeftMargin{14. }%
\CSLRightInline{{Vohra S, Shamseer L, Sampson M, et al.} {CONSORT} extension for reporting {N}-of-1 trials ({CENT}) 2015 statement. \emph{Journal of Clinical Epidemiology}. 2015;76:9-17. doi:\href{https://doi.org/10.1016/j.jclinepi.2015.05.004}{10.1016/j.jclinepi.2015.05.004}}

\bibitem[\citeproctext]{ref-shamseer2015cent}
\CSLLeftMargin{15. }%
\CSLRightInline{{Shamseer L, Sampson M, Bukutu C, Schmid CH, et al.} {CENT} 2015: Explanation and elaboration. \emph{Journal of Clinical Epidemiology}. 2015;76:18-46. doi:\href{https://doi.org/10.1016/j.jclinepi.2015.05.018}{10.1016/j.jclinepi.2015.05.018}}

\bibitem[\citeproctext]{ref-kwasnicka2019}
\CSLLeftMargin{16. }%
\CSLRightInline{{Kwasnicka D, Inauen J, Nieuwenboom W, et al.} Challenges and solutions for {N}-of-1 design studies in health psychology. \emph{Health Psychology Review}. 2019;13(2):163-178. doi:\href{https://doi.org/10.1080/17437199.2018.1564627}{10.1080/17437199.2018.1564627}}

\bibitem[\citeproctext]{ref-mallinckrodt2003}
\CSLLeftMargin{17. }%
\CSLRightInline{Mallinckrodt CH, Clark SWS, Carroll RJ, Molenberghs G. Assessing response profiles from incomplete longitudinal clinical trial data under regulatory considerations. \emph{Journal of Biopharmaceutical Statistics}. 2003;13(2):179-190. doi:\href{https://doi.org/10.1081/BIP-120019265}{10.1081/BIP-120019265}}

\bibitem[\citeproctext]{ref-mallinckrodt2001dropout}
\CSLLeftMargin{18. }%
\CSLRightInline{Mallinckrodt CH, Clark WS, David SR. Accounting for dropout bias using mixed-effects models. \emph{Journal of Biopharmaceutical Statistics}. 2001;11(1-2):9-21. doi:\href{https://doi.org/10.1081/BIP-100104194}{10.1081/BIP-100104194}}

\bibitem[\citeproctext]{ref-mallinckrodt2004duloxetine}
\CSLLeftMargin{19. }%
\CSLRightInline{Mallinckrodt CH, Raskin J, Wohlreich MM, Watkin JG, Detke MJ. The efficacy of duloxetine: A comprehensive summary of results from {MMRM} and {LOCF\_ANCOVA} in eight clinical trials. \emph{BMC Psychiatry}. 2004;4:26. doi:\href{https://doi.org/10.1186/1471-244X-4-26}{10.1186/1471-244X-4-26}}

\bibitem[\citeproctext]{ref-prakash2008}
\CSLLeftMargin{20. }%
\CSLRightInline{Prakash A, Risser RC, Mallinckrodt CH. The impact of analytic method on interpretation of outcomes in longitudinal clinical trials. \emph{International Journal of Clinical Practice}. 2008;62(8):1147-1158. doi:\href{https://doi.org/10.1111/j.1742-1241.2008.01808.x}{10.1111/j.1742-1241.2008.01808.x}}

\bibitem[\citeproctext]{ref-liang1986}
\CSLLeftMargin{21. }%
\CSLRightInline{Liang KY, Zeger SL. Longitudinal data analysis using generalized linear models. \emph{Biometrika}. 1986;73(1):13-22. doi:\href{https://doi.org/10.1093/biomet/73.1.13}{10.1093/biomet/73.1.13}}

\bibitem[\citeproctext]{ref-wang2016geesmv}
\CSLLeftMargin{22. }%
\CSLRightInline{Wang M, Kong L, Li Z, Zhang L. Covariance estimators for generalized estimating equations ({GEE}) in longitudinal analysis with small samples. \emph{Statistics in Medicine}. 2016;35(10):1706-1721. doi:\href{https://doi.org/10.1002/sim.6817}{10.1002/sim.6817}}

\bibitem[\citeproctext]{ref-pan2001}
\CSLLeftMargin{23. }%
\CSLRightInline{Pan W. Sample size and power calculations with correlated binary data. \emph{Controlled Clinical Trials}. 2001;22(3):211-227. doi:\href{https://doi.org/10.1016/s0197-2456(01)00131-3}{10.1016/s0197-2456(01)00131-3}}

\bibitem[\citeproctext]{ref-tong2024stepwedge}
\CSLLeftMargin{24. }%
\CSLRightInline{{Tong G, Nevins P, Ryan M, others, Taljaard M}. A review of current practice in the design and analysis of extremely small stepped-wedge cluster randomized trials. \emph{Clinical Trials}. 2024;22(1):45-56. doi:\href{https://doi.org/10.1177/17407745241276137}{10.1177/17407745241276137}}

\bibitem[\citeproctext]{ref-hedges2022power}
\CSLLeftMargin{25. }%
\CSLRightInline{Hedges LV, Shadish WR, Natesan Batley P. Power analysis for single-case designs: Computations for {\((AB)^k\)} designs. \emph{Behavior Research Methods}. 2023;55(7):3494-3503. doi:\href{https://doi.org/10.3758/s13428-022-01971-9}{10.3758/s13428-022-01971-9}}

\bibitem[\citeproctext]{ref-bouwmeester2020}
\CSLLeftMargin{26. }%
\CSLRightInline{Bouwmeester S, Jongerling J. Power of a randomization test in a single case multiple baseline {AB} design. \emph{PLoS ONE}. 2020;15(2):e0228355. doi:\href{https://doi.org/10.1371/journal.pone.0228355}{10.1371/journal.pone.0228355}}

\bibitem[\citeproctext]{ref-schork2023precision}
\CSLLeftMargin{27. }%
\CSLRightInline{Schork NJ, Beaulieu-Jones B, Liang WS, Smalley S, Goetz LH. Exploring human biology with {N}-of-1 clinical trials. \emph{Cambridge Prisms: Precision Medicine}. 2023;1:e12. doi:\href{https://doi.org/10.1017/pcm.2022.15}{10.1017/pcm.2022.15}}

\bibitem[\citeproctext]{ref-zucker2010combining}
\CSLLeftMargin{28. }%
\CSLRightInline{Zucker DR, Ruthazer R, Schmid CH. Individual ({N}-of-1) trials can be combined to give population comparative treatment effect estimates: Methodologic considerations. \emph{Journal of Clinical Epidemiology}. 2010;63(12):1312-1323. doi:\href{https://doi.org/10.1016/j.jclinepi.2010.04.020}{10.1016/j.jclinepi.2010.04.020}}

\bibitem[\citeproctext]{ref-defelippe2023}
\CSLLeftMargin{29. }%
\CSLRightInline{{Defelippe VM, Thiel GJMW van, Otte WM, Schutgens REG, Stunnenberg B, et al.} Toward responsible clinical {n}-of-1 strategies for rare diseases. \emph{Drug Discovery Today}. 2023;28(10):103688. doi:\href{https://doi.org/10.1016/j.drudis.2023.103688}{10.1016/j.drudis.2023.103688}}

\bibitem[\citeproctext]{ref-berger2003runin}
\CSLLeftMargin{30. }%
\CSLRightInline{Berger VW, Rezvani A, Makarewicz VA. Direct effect on validity of response run-in selection in clinical trials. \emph{Controlled Clinical Trials}. 2003;24(2):156-166. doi:\href{https://doi.org/10.1016/s0197-2456(02)00316-1}{10.1016/s0197-2456(02)00316-1}}

\bibitem[\citeproctext]{ref-senn2002crossover}
\CSLLeftMargin{31. }%
\CSLRightInline{Senn S. \emph{Cross-over Trials in Clinical Research}. 2nd ed. Wiley; 2002.}

\bibitem[\citeproctext]{ref-maxwelldelaney2018}
\CSLLeftMargin{32. }%
\CSLRightInline{Maxwell SE, Delaney HD, Kelley K. \emph{Designing Experiments and Analyzing Data: A Model Comparison Perspective}. 3rd ed. Routledge; 2018.}

\bibitem[\citeproctext]{ref-mallinckrodt2008recommendations}
\CSLLeftMargin{33. }%
\CSLRightInline{Mallinckrodt CH, Lane PW, Schnell D, Peng Y, Mancuso JP. Recommendations for the primary analysis of continuous endpoints in longitudinal clinical trials. \emph{Drug Information Journal}. 2008;42(4):303-319. doi:\href{https://doi.org/10.1177/009286150804200402}{10.1177/009286150804200402}}

\bibitem[\citeproctext]{ref-mauchly1940}
\CSLLeftMargin{34. }%
\CSLRightInline{Mauchly JW. Significance test for sphericity of a normal n-variate distribution. \emph{Annals of Mathematical Statistics}. 1940;11(2):204-209. doi:\href{https://doi.org/10.1214/aoms/1177731915}{10.1214/aoms/1177731915}}

\bibitem[\citeproctext]{ref-mancl2001covariance}
\CSLLeftMargin{35. }%
\CSLRightInline{Mancl LA, DeRouen TA. A covariance estimator for {GEE} with improved small-sample properties. \emph{Biometrics}. 2001;57(1):126-134.}

\end{CSLReferences}

\end{document}
